\documentclass{article}
\usepackage{graphicx} % Required for inserting images
\usepackage{amsmath}
\usepackage{multirow}
\usepackage{hyperref}
\usepackage{xcolor}
\usepackage{tcolorbox}

\definecolor{light}{HTML}{EDF2F7}
\newtcolorbox{keyidea}{
    colback=light,
    colframe=light,
    boxrule=0pt,
    arc=2pt,
    left=9pt, right=9pt,
    top=7pt, bottom=7pt
}


\title{RGA Krypton Measurement Uncertainties and Sensitivity}
\author{Dacheng Xu}
\date{November 2026}

\begin{document}

\maketitle

\section{Uncertainty of Partial Pressure}

The pressure-versus-time data collected when the RGA is exposed only to the ice primer use the subscript $_\mathrm{bkg}$, while the data collected when the MFC is running and xenon is freezing use the subscript $_\mathrm{flw}$. The assumptions are shown in blue boxes.

\begin{keyidea}
    The time intervals between successive groups of AMU measurements are equal.
\end{keyidea}

The background and flow runs contain $N_\mathrm{bkg}$ and $N_\mathrm{flw}$ measurements, respectively. Let $\Delta p_k(A) = p_k(A) - p_k(87)$ be the reference-subtracted partial pressure for measurement group $k$.

The average reference-subtracted partial pressures are
\begin{align}
    P_\mathrm{flw}(A) &= \frac{1}{N_\mathrm{flw}}\sum_{i=0}^{N_\mathrm{flw}-1}\Delta p_i(A), \\
    P_\mathrm{bkg}(A) &= \frac{1}{N_\mathrm{bkg}}\sum_{j=0}^{N_\mathrm{bkg}-1}\Delta p_j(A).
\end{align}

We define the sample partial pressure as
\begin{align}
    P_\mathrm{sample}(A) &= P_\mathrm{flw}(A) - P_\mathrm{bkg}(A).
\end{align}
We assume that $P_\mathrm{sample}(A)$ is a linear function of the true concentration.

\begin{keyidea}
    Let $\sigma_k(A)$ be the standard deviation of the raw reading $p_k(A)$ and $\sigma_0$ that of the raw reference reading $p_k(87)$. We model $\sigma_k(A)$ as an affine function of $\Delta p_k(A)$.
    \begin{align}
        \sigma_k(A) &= \sigma_0\left[1 + c \Delta p_k(A)\right].
    \end{align}
    The reference-channel standard deviation is constant within each run, with values $\sigma_{0,\mathrm{flw}}$ and $\sigma_{0,\mathrm{bkg}}$. We assume a common coefficient $c$ for both runs and estimate it from the background run.
    We assume that the two channel readings in each group are uncorrelated and that readings from distinct groups are uncorrelated, including across runs.
\end{keyidea}

The variances of $P_\mathrm{flw}(A)$ and $P_\mathrm{bkg}(A)$ are
\begin{align}
    \sigma^2_{P_\mathrm{flw}(A)} &= \frac{1}{N_\mathrm{flw}^2}\sum_{i=0}^{N_\mathrm{flw}-1}
        \sigma_i^2(A) + \frac{\sigma^2_{0,\mathrm{flw}}}{N_\mathrm{flw}}, \\
    \sigma^2_{P_\mathrm{bkg}(A)} &= \frac{1}{N_\mathrm{bkg}^2}\sum_{j=0}^{N_\mathrm{bkg}-1}
        \sigma_j^2(A) + \frac{\sigma^2_{0,\mathrm{bkg}}}{N_\mathrm{bkg}}.
\end{align}

The variance of $P_\mathrm{sample}(A)$ is
\begin{align}
    \sigma^2_{P_\mathrm{sample}(A)} &= \sigma^2_{P_\mathrm{flw}(A)} + \sigma^2_{P_\mathrm{bkg}(A)}.
\end{align}

\begin{keyidea}
    During the background run, assume that $\Delta p_j(A)$ stays close to $P_\mathrm{bkg}(A)$, so
    \begin{align}
        \sigma_j(A) &\approx \sigma_{0,\mathrm{bkg}}\left[1 + c P_\mathrm{bkg}(A)\right].
    \end{align}
\end{keyidea}

The background contribution is then
\begin{align}
    \sigma^2_{P_\mathrm{bkg}(A)} &\approx \frac{\sigma^2_{0,\mathrm{bkg}}}{N_\mathrm{bkg}}
        \left\{1 + \left[1 + c P_\mathrm{bkg}(A)\right]^2\right\}.
\end{align}

The flow contribution can be approximated in two limits.

\subsection{Low Concentration}

\begin{keyidea}
    In the low-concentration limit, assume that $\lvert c\Delta p_i(A)\rvert \ll 1$ throughout the flow run.
    \begin{align}
        \sigma_i^2(A) &= \sigma^2_{0,\mathrm{flw}}\left[1 + c\Delta p_i(A)\right]^2 \\
        &\approx \sigma^2_{0,\mathrm{flw}}\left[1 + 2c\Delta p_i(A)\right].
    \end{align}
\end{keyidea}

Thus,
\begin{align}
    \sigma^2_{P_\mathrm{sample}(A)} &\approx \frac{\sigma^2_{0,\mathrm{flw}}}{N_\mathrm{flw}}
        \left[2 + 2cP_\mathrm{flw}(A)\right] + \sigma^2_{P_\mathrm{bkg}(A)}.
\end{align}

\subsection{High Concentration}

\begin{keyidea}
    In the high-concentration limit, assume that $c\Delta p_i(A) \gg 1$ throughout the flow run.
    \begin{align}
        \sigma_i^2(A) &= \sigma^2_{0,\mathrm{flw}}\left[1 + c\Delta p_i(A)\right]^2 \\
        &\approx \sigma^2_{0,\mathrm{flw}} c^2\Delta p_i(A)^2.
    \end{align}
\end{keyidea}

In this limit, the flow-run reference-channel variance is negligible compared with the pressure-dependent term, so
\begin{align}
    \sigma^2_{P_\mathrm{sample}(A)} &\approx \frac{\sigma^2_{0,\mathrm{flw}}c^2}{N_\mathrm{flw}^2}
        \sum_{i=0}^{N_\mathrm{flw}-1}\Delta p_i(A)^2 + \sigma^2_{P_\mathrm{bkg}(A)}.
\end{align}
If $\Delta p_i(A)\approx P_\mathrm{flw}(A)$ throughout the flow run, this simplifies to
\begin{align}
    \sigma^2_{P_\mathrm{sample}(A)} &\approx \frac{\sigma^2_{0,\mathrm{flw}}c^2P^2_\mathrm{flw}(A)}{N_\mathrm{flw}}
        + \sigma^2_{P_\mathrm{bkg}(A)}.
\end{align}

We estimate $\sigma_{0,\mathrm{flw}}$ and $\sigma_{0,\mathrm{bkg}}$
using the sample standard deviations of $p_i(87)$ and $p_j(87)$,
respectively. Let $\sigma_\mathrm{bkg}(A)$ denote the sample standard
deviation of the raw background readings $p_j(A)$. For
$P_\mathrm{bkg}(A) \ne 0$, we estimate $c$ as
\begin{align}
    c &= \frac{\sigma_\mathrm{bkg}(A)/\sigma_{0,\mathrm{bkg}} - 1}{P_\mathrm{bkg}(A)}.
\end{align}

\end{document}
